\chapter{Planejamento do Estudo de Caso}
\label{ch:proposta}

Inicialmente, tendo como base a ISO 25010, utilizando a característica de manutenabilidade com sua subcaracterística de testabilidade, utilizaríamos métricas que empregavam similaridade semântica e similaridade de cosseno para responder às questões específicas: “Em que medida a clareza das sentenças BDD influencia a capacidade da IA de identificar corretamente elementos da interface gráfica e gerar testes alinhados ao comportamento descrito no cenário?” e “Em que medida a padronização da escrita dos cenários BDD contribui para a testabilidade ao melhorar a coerência e a aderência semântica do cenário durante a validação semântica prévia à geração do código, reduzindo ambiguidades e facilitando a correta compreensão do comportamento pela IA?”.

Elas seriam utilizadas na solução que criamos, a qual havia sido elaborada a partir do conhecimento obtido nos artigos selecionados após a revisão estruturada. Ela era dividida em módulos, nos quais o Módulo 1 recebe as histórias de usuário e suas versões em BDD, processa linguisticamente os cenários com técnicas como tokenização, POS tagging e análise de papéis semânticos, e gera um JSON estruturado com os elementos extraídos dos requisitos. O Módulo 2 valida semanticamente se o JSON preserva o significado do cenário BDD original por meio de embeddings e similaridade por cosseno; se a qualidade for suficiente, recupera contexto com RAG e gera automaticamente os testes em Cypress com apoio de um LLM. O Módulo 3 compara semanticamente os requisitos com os testes gerados para identificar inconsistências e utiliza um agente de IA para revisar e ajustar o código antes da comparação final com testes produzidos por modelos generalistas.

Porém, devido à dificuldade de medir similaridade semântica com precisão, abandonamos essas questões específicas e fomos para uma outra abordagem, visando medir eficácia e eficiência, utilizando duas questões específicas: “Em que medida a estruturação e o detalhamento dos cenários BDD influenciam a eficácia da IA na geração de testes automatizados?” e “Em que medida a estruturação e o detalhamento dos cenários BDD influenciam a eficiência do processo de geração automática de testes, reduzindo a intervenção manual pós-geração, especialmente em termos de correção, complementação de contexto, regeneração e descarte dos artefatos gerados, bem como o tempo e o custo do processo?”.

Mas, devido à dificuldade de encontrar métricas que pudessem responder à segunda questão específica, medindo a eficiência, resolvemos voltar aos artigos da revisão estruturada e reler os artigos para encontrar as métricas. Porém, apesar de não conseguir achar as métricas que queríamos, encontramos, no artigo Cypress Copilot: Development of an AI Assistant for Boosting Productivity and Transforming Web Application Testing, um caminho diferente, já que sua solução, que é uma extensão do Visual Studio Code, na qual o usuário insere o BDD e a API key e escolhe o modelo de IA, como GPT-4, gera o teste automatizado em Cypress e o Page Object Model do teste automatizado, o que é muito próximo do que precisamos, mas, como é de código aberto, podemos alterá-la para o nosso caso.

O próprio artigo oferece métricas que medem eficácia e qualidade do código gerado, que iremos utilizar. A questão de pesquisa e o objetivo central do TCC permaneceram inalterados, e essa mudança ocorreu antes da execução e implementação do estudo.

\input{editaveis/secao-proposta/definicao}